% ============================================================
% ch00.tex —— 序章：如何烧掉一本数学书
% ============================================================
\chapter*{序章\quad 如何烧掉一本数学书}
\addcontentsline{toc}{chapter}{序章\quad 如何烧掉一本数学书}

这本书的第一件事，是把你已经"学会"的东西全部没收。

别紧张，不是永久没收。只是接下来三百页里，我们要把数学当成一间毛坯房：不许住现成的，全部自己装修。作为热身，这一章我们做三件事：\textbf{发明一个运算、发明一个数、发明一条规则}。你会发现"发明数学"这门手艺，门槛低得出奇，玩起来却深不见底。

\section{第一次发明：让两个数"相爱"}

\begin{kaiwei}
自己发明一种运算：规定 $a \,\heartsuit\, b = a + b + 1$。比如 $2\,\heartsuit\,3 = 2+3+1 = 6$，$10\,\heartsuit\,(-2) = 10-2+1 = 9$。
\begin{enumerate}[label=(\arabic*)]
  \item 它遵守交换律吗？（$a\,\heartsuit\,b$ 和 $b\,\heartsuit\,a$ 一样吗？）
  \item 有没有一个数 $e$，使得任何数 $a$ 满足 $a \,\heartsuit\, e = a$？（一个"什么都不做"的数）
  \item 每个数 $a$ 有没有一个"搭档" $b$，使得 $a \,\heartsuit\, b$ 恰好等于那个"什么都不做"的数？
\end{enumerate}
\end{kaiwei}

先自己玩五分钟，拿起笔真的算几个，再往下读。\textbf{本书的规矩从现在开始生效：每个"开胃问题"都值得你先动手，哪怕只算两分钟}——因为你亲手算过的东西，后面读到时会长出记忆；直接读答案的东西，明天就忘。

(1) \textbf{交换律}。按定义展开两边：
\[
a\,\heartsuit\,b = a+b+1, \qquad b\,\heartsuit\,a = b+a+1 .
\]
普通加法有交换律，$a+b = b+a$，所以两个结果相等。$\heartsuit$ 遵守交换律。\checkmark

(2) \textbf{"什么都不做"的数}。我们要找一个 $e$，让\textbf{任何} $a$ 都满足 $a \,\heartsuit\, e = a$。把左边按定义展开：
\[
a \,\heartsuit\, e = a + e + 1 .
\]
要它等于 $a$，就得 $e + 1 = 0$，解得
\[
\boxed{e = -1} .
\]
验证三组：$5\,\heartsuit\,(-1) = 5 - 1 + 1 = 5$；$100\,\heartsuit\,(-1) = 100$；$(-7)\,\heartsuit\,(-1) = -7$。果然什么都不做！在普通加法里这个角色由 $0$ 扮演，在"$\heartsuit$"的世界里，\Keyword{戏份被 $-1$ 抢走了}。

这里值得停一秒体会：我们问的问题——"有没有一个什么都不做的数"——对\textbf{任何}运算都可以问。它不是关于 $\heartsuit$ 的怪问题，而是关于"运算"这个概念本身的普遍问题。问出普遍的问题，是数学家的第一个标志。

(3) \textbf{搭档}。$3$ 的搭档是谁？按定义，$3 \,\heartsuit\, b = 3 + b + 1 = b+4$，我们要求它等于"什么都不做"的数 $-1$：
\[
b + 4 = -1 \quad\Longrightarrow\quad b = -5 .
\]
验证：$3\,\heartsuit\,(-5) = 3 - 5 + 1 = -1$。\checkmark 一般地，$a$ 的搭档是 $-a-2$（解方程 $a + b + 1 = -1$）。

你刚刚做完的这场游戏，在大学里有一个吓人的名字，叫\Keyword{验证群的公理}（第 7 章）。看，你还没上大学就已经会用它的全部内容了，缺的只是一张正式发票。这就是本书的做法：\textbf{先玩，后命名；先发明，后盖章}。

\section{第二次发明：一个被骂了三百年的数}

今天你觉得负数天经地义：零下 $3^\circ$C、欠债 $5$ 元、海拔 $-155$ 米的死海、电梯负一层。可是在欧洲，直到十六、十七世纪，很多数学家还在管负数叫"\Keyword{比没有还少的荒谬之数}"（numerus absurdus，荒谬数）。帕斯卡尔——对，就是那个大名鼎鼎的帕斯卡尔——说"从 $0$ 减去 $4$，纯粹是胡说八道"。他们不是笨，他们只是死守一条规矩：数是用来\textbf{数东西}的，而世界上没有"$-3$ 个苹果"。

\textbf{负数是怎么打赢这场舆论战的？}不是靠嘴仗，是靠"有用"。两件武器：

\textbf{武器一：记账。}中世纪欧洲的商船会计发明了复式记账：资产记一边、负债记另一边。负债本质上就是"负的资产"——$-500$ 元的负债加 $500$ 元的现金，合计正好 $0$。商人们不管哲学家怎么骂，账本能平就是硬道理。负数在账本的掩护下率先入伍。

\textbf{武器二：解方程。}考虑方程 $x + 5 = 3$。守旧派说"无解，因为 $3$ 比 $5$ 小"；接受负数的人说"解是 $-2$，验算：$-2 + 5 = 3$，成立"。更妙的是，很多\textbf{中间步骤里出现负数、最后答案却是正数}的题目，只有允许负数才能算到底。比如按古代的流程解某些问题，中途不得不写"$5 - 12$"——守旧派在这里就卡死了，而"荒谬派"继续前进，最后落回一个完全合理的正数答案。\textbf{允许负数，不产生任何错误答案，只解锁了大量算不动的题}——这笔账划得来。

于是十七世纪之后，负数悄悄转正。注意这段历史里的发明套路，它将贯穿全书：
\begin{jielun}
旧工具遇到了做不了的事 $\Longrightarrow$ \textbf{发一个新工具} $\Longrightarrow$ 规定新工具怎么用（不许自相矛盾） $\Longrightarrow$ 检查：旧麻烦解决了吗？新麻烦出现了吗？
\end{jielun}

\textbf{分数}走的是同一条路。四千年前的埃及人只有\textbf{单位分数}（分子是 $1$ 的分数：$\tfrac12, \tfrac13, \tfrac15,\dots$），兰德纸草书（约公元前 1650 年）里有一张表，专门教人把 $\tfrac25, \tfrac27, \tfrac29$ 这类分数拆成单位分数之和，比如
\[
\frac{2}{7} = \frac{1}{4} + \frac{1}{28}, \qquad \frac{2}{13} = \frac{1}{8} + \frac{1}{52} + \frac{1}{104}.
\]
为什么这么费劲？因为通用的分数还没被发明，他们只有"几份中的一份"这个直觉。后人发明了"分子随便取"的分数，除法才全面畅通。第 10 章你会亲眼看到一台\textbf{把整数加工成分数的完整流水线}——不是比喻，是真正的施工图。

\textbf{零}的故事更曲折。罗马数字里根本没有 $0$：一、五、十、五十……每个符号都代表"实实在在的东西"，"什么都没有"凭什么要一个符号？麻烦出在哪？你用罗马数字算一道乘法就懂了——没有位值，就没有速算。而位值记数法（个位、十位、百位……）有一个致命的漏洞：$1$ 后面空一格表示"十"，可空格怎么和"这里什么都没有"区分？$205$ 和 $25$ 差了两百多，写在纸上却只差一个空白。印度的数学家（约公元 5 世纪）给这个空白发明了一个符号——$0$，"无"第一次有了名字。这个符号随阿拉伯人的商路传遍世界，所以今天我们管它叫"阿拉伯数字"。顺便说，零还有一层身份：它是\textbf{加法里"什么都不做"的数}——你在序章第一个问题里刚刚亲手为 $\heartsuit$ 世界找过同类角色（那里是 $-1$）。同一个点子，印度人用了一千年，你用了一道习题的时间。

\section{第三次发明：一条规则}

最后热身一个"发明规则"的例子，它会在第 5 章大放异彩。

\begin{kaiwei}
今天是星期六。问：$100$ 天后是星期几？
\end{kaiwei}

没人真的去数 $100$ 天。聪明的做法分三步：

\textbf{第一步：找循环。}星期几以 $7$ 天为周期循环——过了 $7$ 天回到原样，过了 $14$ 天还是原样，过了 $7$ 的任何倍数天依然原样。所以"$100$ 天后"和"$100$ 去掉所有整星期之后剩的那几天后"是同一件事。

\textbf{第二步：做带余除法。}
\[
100 = 14 \times 7 + 2 .
\]
$100$ 天里藏着 $14$ 个完整的星期，剩下零头 $2$ 天。

\textbf{第三步：只看零头。}星期六再过 $2$ 天：星期日（$1$ 天后）、\textbf{星期一}（$2$ 天后）。答案：星期一。

看清楚我们刚才干了什么：$100$ 这个数被"压扁"成了 $2$，而所有关于星期几的问题，答案在压扁之后\textbf{完全不变}。起作用的不是 $100$ 本身，而是它除以 $7$ 的\textbf{余数}。这一压，压出了一整个新世界——\Keyword{余数的世界}：一个只有 $7$ 个数的小王国，加法、乘法一应俱全，第 1 章我们就搬进去住。

再预告一下这个小王国的威力。两千年前，《孙子算经》用同样的"压扁术" solved 一个"物不知数"的谜题（第 5 章）；今天，你每一笔网络支付背后的密码锁，锁芯就是余数（第 6 章）。余数是数学史上性价比最高的发明之一：工具成本几乎为零（就是除法的零头），从军营点兵一路用到国际结算。

\section{本书使用说明书}

\begin{itemize}
  \item \textbf{土名与正名}：每个术语出场前都先干脏活。见到"落脚点""瞬间变化率""堆积"这些土名时放心用；它们的正式户口（极限、导数、积分）会在你用顺手之后才发给你。附录 A 有完整的对照表——读完这本书，你可以拿它去大学课本里"对暗号"。
  \item \textbf{三条阅读路线}：
  \begin{itemize}
    \item \textbf{速通线}（想快速尝到"高数"味道）：序章 $\to$ 第 1、5、7、15、17、23、27、32 章 $\to$ 尾声；
    \item \textbf{数论线}（喜欢整数与密码谜题）：篇一 $\to$ 第 7--11 章 $\to$ 篇四；
    \item \textbf{几何线}（喜欢图形与空间）：篇三 $\to$ 篇五 $\to$ 篇六，再回头看第 8 章。
  \end{itemize}
  六篇之间的依赖关系画成地图，就是这样：

\begin{center}
\begin{tikzpicture}[node distance=0.9cm, every node/.style={font=\small}]
  \node[draw=shenlan!70, rounded corners, fill=shenlan!8] (p1) at (0,0) {篇一 数论};
  \node[draw=shenlan!70, rounded corners, fill=shenlan!8] (p2) at (0,-1.4) {篇二 抽象代数};
  \node[draw=cheng!80, rounded corners, fill=cheng!8] (p3) at (4.2,0) {篇三 分析};
  \node[draw=zi!70, rounded corners, fill=zi!8] (p4) at (2.1,-2.8) {篇四 解析数论};
  \node[draw=molv!70, rounded corners, fill=molv!8] (p5) at (4.2,-1.4) {篇五 微分几何};
  \node[draw=shaohong!70, rounded corners, fill=shaohong!8] (p6) at (2.1,-4.2) {篇六 代数几何};
  \draw[-{Stealth}] (p1) -- (p2);
  \draw[-{Stealth}] (p1) -- (p4);
  \draw[-{Stealth}] (p3) -- (p4);
  \draw[-{Stealth}] (p3) -- (p5);
  \draw[-{Stealth}] (p2) -- (p6);
  \draw[-{Stealth}] (p5) -- (p6);
  \draw[-{Stealth}] (p3) to[bend left=18] (p6);
\end{tikzpicture}
\end{center}

  箭头的意思是"箭头尾的篇为箭头头的篇提供弹药"。你可以看到：篇三（分析）的弹药送得最远——它同时供给篇四、五、六；而篇一和篇二手拉手，一起撑起篇六的终章。
  \item \textbf{星号规则}：$\star$ 动脑不动笔，$\star\star$ 纸笔十分钟能搞定，$\bigstar\bigstar\bigstar$ 选做，做不出来不丢人——有几道题全人类（包括职业数学家）都还没做出来，第 24 章你会见到一道价值一百万美元的。
  \item \textbf{诚实声明}：凡是我说"我们跳过了证明"的地方，都是真的跳过了，不是"留给读者作为练习"的客套话。地毯下面有什么，我在声明里写清楚，并告诉你哪本大学教材里有完整版本（附录 D）。这本书不装神秘。
\end{itemize}

\begin{dongshou}
\begin{itemize}
  \yisi 证明"$\heartsuit$"运算还满足结合律：$(a\,\heartsuit\,b)\,\heartsuit\,c = a\,\heartsuit\,(b\,\heartsuit\,c)$。（把两边都按定义展开：左边 $= (a+b+1)+c+1$，右边 $= a+(b+c+1)+1$，都等于 $a+b+c+2$。注意"多加的那个 $1$"两边各出现两次——数清楚它，就是证明的全部。）
  \yier 再发明一个你自己的运算 $a\,\clubsuit\,b = a + 2b$。它有交换律吗？有"什么都不做"的数吗？每个数都有搭档吗？（提示：解方程 $a + 2e = a$ 得 $e = 0$，看起来不错；但找 $3$ 的搭档要解 $3 + 2b = 0$，得 $b = -\tfrac32$——\textbf{不是整数}。也就是说在"整数俱乐部"里 $\clubsuit$ 有人没有搭档；把俱乐部扩招成"分数世界"，搭档就齐了。这个"换个会员范围让体检过关"的思路，第 7、10 章还要反复用。）
  \yier 你肯定见过 $a \times 0 = 0$。请模仿本章的精神追问一句：\textbf{为什么}任何数乘 $0$ 都得 $0$？试着只用两条已知的规则——分配律 $a\times(b+c) = a\times b + a\times c$，以及"$0$ 是加法里什么都不做的数"——把它证明出来。（提示：$0 = 0 + 0$，两边同乘 $a$，得 $a\times 0 = a\times 0 + a\times 0$，然后在等式两边同时"撤走"一个 $a\times 0$。你在做的事情，大学里叫"用公理证明 $0$ 是吸收元"。）
\end{itemize}
\end{dongshou}

\begin{shaonao}
把"$\heartsuit$"继续玩下去：$a \,\heartsuit\, a = 2a + 1$；$a \,\heartsuit\, a \,\heartsuit\, a$（约定从左往右算）$= (2a+1) + a + 1 = 3a + 2$；四个 $a$ 连做得 $4a+3$。看出规律了吗？$n$ 个 $a$ 连做等于 $na + (n-1)$。现在做一次"换算"：令 $u = a + 1$，你会发现 $na + (n-1) + 1 = n(a+1) = nu$——也就是说，把每个数都加 $1$ 换算过去，$\heartsuit$ 运算\textbf{整个变成了普通加法}！这就是第 7 章要讲的"换一套单位元，世界照常运转"，数学家管它叫\textbf{同构}（第 8 章）：你发明的运算，骨子里就是加法穿了件马甲。
\end{shaonao}

\begin{chengshi}
本书的"发明顺序"是为\textbf{理解}设计的最短路线，不是历史原路。比如历史上分数远早于负数的普及（埃及人用单位分数时，欧洲还在为负数吵架）；"物不知数"的口诀比中国剩余定理的一般形式早了一千多年。真实的数学史绕了远路、吵了架、死了人（第 12 章会讲一位真实的牺牲者）。想看历史的原貌，见附录 D 的书单。
\end{chengshi}

热身完毕。下一章，我们收拾行李，搬进一个只有七个数的小王国。
